\section*{Chapitre \ref{ch:g10:chemical-species} --- Espèces chimiques naturelles et synthétiques}

\begin{solution}{exo:g10:chemical-species:1}
Naturelles~: la vanilline d'une gousse, la caféine des grains de café.
Synthétiques~: l'aspirine, la vanilline d'usine.
\end{solution}

\begin{solution}{exo:g10:chemical-species:2}
Ce sont la même \omterm{def:g6:pure-substances-and-mixtures:species}{espèce chimique}, et une espèce a les mêmes propriétés
quelle que soit son origine.
\end{solution}

\begin{solution}{exo:g10:chemical-species:3}
$R_f = 2.4/6.0 = 0.40$.
\end{solution}

\begin{solution}{exo:g10:chemical-species:4}
Il refroidit la vapeur et la retransforme en liquide. En entrant par le
bas, l'eau remplit toute la double enveloppe (et rencontre en dernier la
vapeur la plus chaude, ce qui refroidit le mieux).
\end{solution}

\begin{solution}{exo:g10:chemical-species:5}
Dissoudre l'espèce beaucoup mieux que l'eau~; être \omterm{def:g6:solutions-and-solubility:miscible}{non miscible} avec
l'eau~; être aussi peu dangereux que possible et facile à éliminer.
\end{solution}

\begin{solution}{exo:g10:chemical-species:6}
L'éthanol est \omterm{def:g6:solutions-and-solubility:miscible}{miscible} avec l'eau~: aucune phase ne se forme, on ne peut
rien séparer. Le cyclohexane n'est pas \omterm{def:g6:solutions-and-solubility:miscible}{miscible} avec l'eau et dissout
bien le linalol.
\end{solution}

\begin{solution}{exo:g10:chemical-species:7}
L'eau est au-dessus~: \qty{1}{mL} de dichlorométhane pèse \qty{1.33}{g},
plus que \qty{1}{mL} d'eau~: il coule.
\end{solution}

\begin{solution}{exo:g10:chemical-species:8}
Elle fond trop bas (aspirine~: \qty{135}{\celsius}) et sur un
intervalle~: l'échantillon est impur, ou ce n'est pas de l'aspirine.
\end{solution}

\begin{solution}{exo:g10:chemical-species:9}
Du linalol et de l'éthanoate de linalyle (ses deux taches sont aux
hauteurs des deux références). L'éthanoate de linalyle est monté le plus
haut.
\end{solution}

\begin{solution}{exo:g10:chemical-species:10}
\omterm{def:g10:chemical-species:distillation}{Hydrodistillation}~; ajout de cyclohexane au distillat et agitation~;
séparation des phases~; \omterm{def:g3:separating-mixtures:evaporate}{évaporation} du cyclohexane.
\end{solution}

\begin{solution}{exo:g10:chemical-species:11}
Le crayon à papier ne \omterm{def:g2:mixing-and-dissolving:dissolve}{se dissout} pas dans l'\omterm{def:g10:chemical-species:tlc}{éluant}~; au-dessus de
l'\omterm{def:g10:chemical-species:tlc}{éluant}, les taches sont entraînées le long de la plaque au lieu de \omterm{def:g2:mixing-and-dissolving:dissolve}{se
dissoudre} dans la cuve.
\end{solution}

\begin{solution}{exo:g10:chemical-species:12}
Le cyclohexane bout à \qty{81}{\celsius}, le linalol à
\qty{198}{\celsius}~: un chauffage doux chasse le cyclohexane tandis que
le linalol reste.
\end{solution}

\begin{solution}{exo:g10:chemical-species:13}
L'échantillon est probablement cette espèce (même $R_f$ que l'\omterm{def:g10:chemical-species:natural}{espèce
naturelle} sur la même plaque) et probablement pur (une seule tache,
fusion nette).
\end{solution}

\begin{solution}{exo:g10:chemical-species:14}
$0.40 \times 5.5 = \qty{2.2}{cm}$. L'autre est à $0.75 \times 5.5 =
\qty{4.1}{cm}$ (au millimètre près), donc environ \qty{1.9}{cm} au-dessus
de A.
\end{solution}

\begin{solution}{exo:g10:chemical-species:15}
\qty{200}{mL} peuvent dissoudre $1.6 \times 0.200 = \qty{0.32}{g}$~: avec
\qty{0.20}{g}, le distillat n'est pas saturé~; il pourrait en dissoudre
encore \qty{0.12}{g}. Les \qty{0.20}{g} dissous seraient perdus avec
l'eau~; le cyclohexane, qui dissout bien mieux le linalol, en récupère la
plus grande partie.
\end{solution}

\begin{solution}{pb:g10:chemical-species:1}
\textbf{1.} Sa vapeur entraîne les espèces odorantes hors des fleurs et
jusque dans le réfrigérant.

\textbf{2.} Il refroidit la vapeur pour la retransformer en liquide~;
l'eau froide entre par l'extrémité inférieure.

\textbf{3.} \omterm{def:g6:solutions-and-solubility:miscible}{Non miscible}~; moins dense que l'eau (elle flotte).

\textbf{4.} La couche est très fine et une partie de l'huile est dissoute
dans l'eau ou dispersée dans celle-ci~: une \omterm{def:g10:chemical-species:extraction}{extraction} permet de la
récupérer.

\textbf{5.} Il dissout bien le linalol, n'est pas \omterm{def:g6:solutions-and-solubility:miscible}{miscible} avec l'eau, et
bout à basse température (\qty{81}{\celsius})~: il est donc facile à
éliminer.

\textbf{6.} Non~: l'éthanol se \omterm{def:g6:pure-substances-and-mixtures:mixture}{mélange} à l'eau et aucune phase ne se
forme.

\textbf{7.} Au-dessus (\qty{0.78}{g} par mL, moins que l'eau). Avec le
dichlorométhane (\qty{1.33}{g} par mL), la phase du \omterm{def:g6:solutions-and-solubility:solute}{solvant} serait
au-dessous.

\textbf{8.} Pas de flamme à proximité (inflammable)~; travailler sous la
hotte avec des gants, et récupérer les déchets au lieu de les verser dans
l'évier (toxique pour les organismes aquatiques).

\textbf{9.} Le cyclohexane bout à \qty{81}{\celsius}, le linalol à
\qty{198}{\celsius}~: un chauffage doux n'élimine que le \omterm{def:g6:solutions-and-solubility:solute}{solvant}.

\textbf{10.} $R_f(\text{L}) = 2.4/6.0 = 0.40$~; $R_f(\text{A}) = 4.5/6.0
= 0.75$.

\textbf{11.} Du linalol et de l'éthanoate de linalyle.

\textbf{12.} Non~: elle contient au moins deux espèces~; c'est un
\omterm{def:g6:pure-substances-and-mixtures:mixture}{mélange}.

\textbf{13.} Pour que toutes les taches soient éluées dans les mêmes
conditions~: c'est seulement ainsi que l'on peut comparer les hauteurs.

\textbf{14.} Son $R_f$ est égal à celui du linalol dans les mêmes
conditions~: c'est probablement du linalol, identique à l'\omterm{def:g10:chemical-species:natural}{espèce
naturelle}.

\textbf{15.} $1.0 \times 0.90 = \qty{0.90}{g}$.

\textbf{16.} $0.90/100 = \qty{0.90}{\percent}$.

\textbf{17.} $100/0.009 \approx 11\,000$\,g, soit environ \qty{11}{kg}
de fleurs.

\textbf{18.} Elle est extraite d'une plante, sans \omterm{def:g5:irreversible-changes:chemical-change}{transformation
chimique}. Mais son linalol est la même espèce que le linalol de
synthèse~: même formule, mêmes propriétés.

\textbf{19.} \qty{0.90}{g} d'huile essentielle à partir de \qty{100}{g}
de fleurs.
\end{solution}
